\chapter{Introduction}
\label{ch:introduction}

\section{Background}

Manufacturing Resource Planning (MRP II) extends basic Material Requirements
Planning. Where MRP answers what material is needed, how much, and when,
MRP II also accounts for production resources: machines, people, finite
capacity, supplier constraints, schedules, costs, and feedback from events on
the shop floor~\cite{siemensmrp2,oracleglossary}.

\todo{expand to roughly two paragraphs on how MRP II is used in practice; cite
a standard reference such as Wight or a current textbook}

\section{Motivation}

Conventional MRP II systems compute a plan centrally. A single scheduler holds
all state and makes all decisions. Two consequences follow. First, a single
disruption tends to invalidate the whole plan, because the plan was produced as
one indivisible calculation. Second, the output is difficult to interrogate: the
plan states what to do, but not which alternatives were considered or why they
lost.

Distributing the decision across independent agents addresses both. Each
resource owns its own constraints and answers for itself, so a failure touches
only the commitments that depended on it. Because a plan is assembled from
individual offers and acceptances, the reasoning is recorded as a by-product of
how the plan was built.

\section{Problem Definition}

Build a decentralized MRP II-lite system in which independent agents negotiate
for inventory, supplier deliveries, and production capacity. The system must
produce a feasible production plan, react to disruptions by renegotiating only
the affected unfinished work, and show why each decision was made.

\section{Objectives}

\begin{enumerate}[nosep]
  \item Explode a multi-level bill of materials into material and operation
        requirements for a production order.
  \item Net requirements against inventory and reserve what is available.
  \item Obtain and compare supplier offers for the resulting shortages.
  \item Schedule operations against finite work-centre capacity.
  \item Negotiate between order agents and resource agents using a simplified
        Contract Net Protocol carried over the A2A protocol.
  \item Convert each accepted offer into exactly one resource commitment.
  \item Renegotiate affected unfinished work after an injected machine failure
        or supplier delay, leaving completed and unaffected commitments fixed.
  \item Record a complete, auditable decision trace and present it in a
        dashboard.
  \item Continue to plan and replan when the optional local language model is
        unavailable.
\end{enumerate}

\section{Scope and Limitations}

The system is an MRP II-lite subset, not a full ERP. It implements BOM
explosion, inventory netting and reservation, supplier selection, finite
capacity scheduling, operation sequencing, cost and due-date trade-offs,
negotiation and commitment, failure-driven replanning, and a decision trace.

The following are explicitly out of scope:

\begin{itemize}[nosep]
  \item Financial accounting, payroll, and purchasing workflows.
  \item Global optimality. The system produces a feasible, explainable plan
        assembled from local decisions; it does not claim an optimal one.
  \item Blockchain or distributed consensus, which the problem does not
        require.
  \item Fully offline operation. The database (Neon) and the primary language
        model provider are cloud services and need outbound internet access.
        Only the optional local model path runs without internet, and only for
        language-model assistance.
\end{itemize}

\section{Organisation of the Report}

\todo{write once the chapters are final -- one sentence per chapter}
